\documentclass[11pt,a4paper]{article}

% --- THIẾT LẬP LỀ TRANG & KHỔ GIẤY ---
\usepackage[top=1.8cm,bottom=1.8cm,left=1.8cm,right=1.8cm]{geometry}

% --- HỖ TRỢ ĐA DẠNG TRÌNH BIÊN DỊCH (XeLaTeX / pdfLaTeX / LuaLaTeX) ---
\usepackage{iftex}
\ifXeTeX
  \usepackage{fontspec}
  \usepackage[vietnamese]{babel}
  \defaultfontfeatures{Ligatures=TeX}
\else
  \usepackage[utf8]{inputenc}
  \usepackage[T5,T1]{fontenc}
  \usepackage[vietnamese]{babel}
  \usepackage{mathptmx}
\fi

% --- CÁC GÓI TOÁN HỌC CHUẨN AMS ---
\usepackage{amsmath}
\usepackage{amssymb}
\usepackage{amsfonts}
\usepackage{mathtools}
\usepackage{bm}

% --- BẢNG BIỂU, DANH SÁCH & TRANG TRÍ ---
\usepackage{tabularx}
\usepackage{booktabs}
\usepackage{array}
\usepackage{enumitem}
\usepackage{multicol}
\usepackage{fancyhdr}
\usepackage{xcolor}

% --- THIẾT LẬP HEADER & FOOTER ---
\pagestyle{fancy}
\fancyhf{}
\renewcommand{\headrulewidth}{0pt}
\lfoot{\textit{\small Bộ đề kiểm tra Toán THCS KNTT - Mã: c1\_6\_bai1\_test3}}
\rfoot{\textit{\small Trang \thepage}}

\begin{document}

% ==============================================================================
% PHẦN ĐẦU ĐỀ THI (HEADER TIÊU CHUẨN BỘ GD&ĐT)
% ==============================================================================
\noindent
\begin{tabular*}{\textwidth}{@{\extracolsep{\fill}}lr@{}}
\textbf{SỞ GD\&ĐT ........................} & \textbf{ĐỀ KIỂM TRA ĐỊNH KỲ NĂM HỌC 2025 - 2026} \\
\textbf{TRƯỜNG THCS ........................} & \textbf{Môn: TOÁN - LỚP 6} \\
\textit{(Đề thi chuẩn ma trận 2025)} & \textit{Thời gian làm bài: 45 phút (không kể giao đề)}
\end{tabular*}

\vspace{6pt}
\begin{center}
{\Large\bfseries Toán 6 (KNTT) - Bài 1: Tập hợp - Bài kiểm tra số 3}\\[4pt]
{\bfseries [ MỨC ĐỘ 3: NÂNG CAO - Mã đề: \texttt{c1\_6\_bai1\_test3} ]}
\end{center}

\vspace{4pt}
\noindent
\rule{\textwidth}{0.8pt}\\[3pt]
Họ và tên thí sinh: \dotfill\ Lớp: \dots\dots\ SBD: \dots\dots\\[-6pt]
\rule{\textwidth}{0.8pt}

\vspace{10pt}

% ==============================================================================
% PHẦN I: TRẮC NGHIỆM NHIỀU PHƯƠNG ÁN (3,0 ĐIỂM)
% ==============================================================================
\noindent
\textbf{\large PHẦN I. Câu trắc nghiệm nhiều phương án lựa chọn (3,0 điểm)}\\
\textit{\small Thí sinh trả lời từ câu 1 đến câu 12. Mỗi câu hỏi thí sinh chỉ chọn một phương án đúng nhất.}
\vspace{6pt}
\noindent\textbf{Câu 1.} [Nhận biết] Cho tập hợp$M = \{x \in \mathbb{N} \mid 5 \le x < 9\}$. Cách viết nào sau đây biểu diễn đúng tập hợp$M$bằng cách liệt kê các phần tử?\par\vspace{2pt}
\noindent\begin{tabular*}{\textwidth}{@{\extracolsep{\fill}}p{0.48\textwidth}p{0.48\textwidth}@{}}
\textbf{A.} $\{5; 6; 7; 8\}$ & \textbf{B.} $\{6; 7; 8; 9\}$ \\
\textbf{C.} $\{5; 6; 7; 8; 9\}$ & \textbf{D.} $\{6; 7; 8\}$
\end{tabular*}\par\vspace{5pt}

\noindent\textbf{Câu 2.} [Nhận biết] Cho tập hợp$A = \{a; b; \{c; d\}; e\}$. Khẳng định nào sau đây là khẳng định ĐÚNG?\par\vspace{2pt}
\noindent\begin{tabular*}{\textwidth}{@{\extracolsep{\fill}}llll@{}}
\textbf{A.} $\{c; d\} \in A$ & \textbf{B.} $c \in A$ & \textbf{C.} $d \in A$ & \textbf{D.} $\{a; b\} \in A$
\end{tabular*}\par\vspace{5pt}

\noindent\textbf{Câu 3.} [Thông hiểu] Cho hai tập hợp$A = \{x \in \mathbb{N} \mid x = 2k + 1,\, k \in \mathbb{N},\, x < 10\}$và$B = \{1; 3; 5; 7; 9; 11\}$. Kết luận nào sau đây đúng?\par\vspace{2pt}
\noindent\begin{tabular*}{\textwidth}{@{\extracolsep{\fill}}p{0.48\textwidth}p{0.48\textwidth}@{}}
\textbf{A.} $B \subset A$ & \textbf{B.} $A$có 6 phần tử \\
\textbf{C.} $A$gồm các phần tử là số lẻ nhỏ hơn$10$ & \textbf{D.} $A = B$
\end{tabular*}\par\vspace{5pt}

\noindent\textbf{Câu 4.} [Thông hiểu] Tập hợp$P$các chữ cái khác nhau trong cụm từ "NGHIEN CUU KHOA HOC" có bao nhiêu phần tử?\par\vspace{2pt}
\noindent\begin{tabular*}{\textwidth}{@{\extracolsep{\fill}}llll@{}}
\textbf{A.} 10 & \textbf{B.} 9 & \textbf{C.} 11 & \textbf{D.} 8
\end{tabular*}\par\vspace{5pt}

\noindent\textbf{Câu 5.} [Thông hiểu] Cho tập hợp$S = \{x \in \mathbb{N} \mid (2x - 6)(x - 12) = 0\}$. Tổng các phần tử của tập hợp$S$bằng:\par\vspace{2pt}
\noindent\begin{tabular*}{\textwidth}{@{\extracolsep{\fill}}llll@{}}
\textbf{A.} 15 & \textbf{B.} 9 & \textbf{C.} 12 & \textbf{D.} 3
\end{tabular*}\par\vspace{5pt}

\noindent\textbf{Câu 6.} [Thông hiểu] Cho tập hợp$X = \{n \in \mathbb{N}^* \mid n < 100 \text{ và } n \text{ chia hết cho cả } 2 \text{ và } 5\}$. Số phần tử của tập hợp$X$là:\par\vspace{2pt}
\noindent\begin{tabular*}{\textwidth}{@{\extracolsep{\fill}}llll@{}}
\textbf{A.} 9 & \textbf{B.} 10 & \textbf{C.} 8 & \textbf{D.} 11
\end{tabular*}\par\vspace{5pt}

\noindent\textbf{Câu 7.} [Thông hiểu] Cho tập hợp$A = \{x \in \mathbb{N} \mid x = 3k + 1,\, k \in \mathbb{N},\, 10 < x < 40\}$. Số phần tử của$A$là:\par\vspace{2pt}
\noindent\begin{tabular*}{\textwidth}{@{\extracolsep{\fill}}llll@{}}
\textbf{A.} 9 & \textbf{B.} 10 & \textbf{C.} 11 & \textbf{D.} 12
\end{tabular*}\par\vspace{5pt}

\noindent\textbf{Câu 8.} [Thông hiểu] Số tập hợp con gồm đúng 2 phần tử của tập hợp$E = \{2; 4; 6; 8\}$là:\par\vspace{2pt}
\noindent\begin{tabular*}{\textwidth}{@{\extracolsep{\fill}}llll@{}}
\textbf{A.} 4 & \textbf{B.} 6 & \textbf{C.} 8 & \textbf{D.} 12
\end{tabular*}\par\vspace{5pt}

\noindent\textbf{Câu 9.} [Vận dụng thấp] Cho tập hợp$A = \{x \in \mathbb{N} \mid \frac{12}{x - 1} \in \mathbb{N}\}$. Số phần tử của tập hợp$A$là:\par\vspace{2pt}
\noindent\begin{tabular*}{\textwidth}{@{\extracolsep{\fill}}llll@{}}
\textbf{A.} 6 & \textbf{B.} 5 & \textbf{C.} 7 & \textbf{D.} 12
\end{tabular*}\par\vspace{5pt}

\noindent\textbf{Câu 10.} [Vận dụng thấp] Cho tập hợp$M = \{x \in \mathbb{N} \mid x = a + b,\, a \in \{1; 3; 5\},\, b \in \{2; 4; 6\}\}$. Số phần tử khác nhau của tập hợp$M$là:\par\vspace{2pt}
\noindent\begin{tabular*}{\textwidth}{@{\extracolsep{\fill}}llll@{}}
\textbf{A.} 5 & \textbf{B.} 7 & \textbf{C.} 9 & \textbf{D.} 6
\end{tabular*}\par\vspace{5pt}

\noindent\textbf{Câu 11.} [Vận dụng thấp] Cho tập hợp$E$có$n$phần tử ($n \in \mathbb{N}^*$). Biết rằng số tập hợp con của$E$gấp 4 lần số tập con của một tập hợp có$3$phần tử. Giá trị của$n$là:\par\vspace{2pt}
\noindent\begin{tabular*}{\textwidth}{@{\extracolsep{\fill}}llll@{}}
\textbf{A.} 5 & \textbf{B.} 6 & \textbf{C.} 7 & \textbf{D.} 8
\end{tabular*}\par\vspace{5pt}

\noindent\textbf{Câu 12.} [Vận dụng thấp] Cho hai tập hợp$A = \{m - 1; 5\}$và$B = \{2; 2m + 1\}$với$m \in \mathbb{N}$. Có bao nhiêu giá trị của$m$để tập hợp$A$và tập hợp$B$có ít nhất một phần tử chung?\par\vspace{2pt}
\noindent\begin{tabular*}{\textwidth}{@{\extracolsep{\fill}}llll@{}}
\textbf{A.} 2 & \textbf{B.} 3 & \textbf{C.} 1 & \textbf{D.} 4
\end{tabular*}\par\vspace{5pt}


% ==============================================================================
% PHẦN II: TRẮC NGHIỆM ĐÚNG / SAI (2,0 ĐIỂM)
% ==============================================================================
\vspace{6pt}
\noindent
\textbf{\large PHẦN II. Câu trắc nghiệm đúng sai (2,0 điểm)}\\
\textit{\small Thí sinh trả lời từ câu 1 đến câu 2. Trong mỗi ý a), b), c), d) ở mỗi câu, thí sinh chọn đúng hoặc sai.}\\
\textit{\small (Quy tắc tính điểm: Đúng 1 ý được 0,1đ; đúng 2 ý được 0,25đ; đúng 3 ý được 0,5đ; đúng 4 ý được 1,0đ).}
\vspace{6pt}
\noindent\textbf{Câu 1.} Cho tập hợp$A = \{x \in \mathbb{N} \mid x = 4k + 1,\, k \in \mathbb{N},\, x \le 45\}$và$B = \{y \in \mathbb{N} \mid y \text{ là ước của } 60\}$.\par\vspace{2pt}
\begin{enumerate}[label=\textbf{\alph*)}, leftmargin=1.2cm, itemsep=2pt, topsep=2pt]
  \item \textnormal{[Thông hiểu]} Số phần tử của tập hợp$A$là$12$.
  \item \textnormal{[Thông hiểu]} Số phần tử của tập hợp$B$là$10$.
  \item \textnormal{[Vận dụng thấp]} Tập hợp các phần tử thuộc cả$A$và$B$có đúng$2$phần tử.
  \item \textnormal{[Vận dụng thấp]} Tổng các phần tử là số lẻ của tập hợp$B$bằng$24$.
\end{enumerate}
\vspace{6pt}

\noindent\textbf{Câu 2.} Cho tập hợp$S = \{1; 2; 3; 4; 5\}$. Xét các tập hợp con của$S$và các điều kiện liên quan đến số phần tử và tổng các phần tử.\par\vspace{2pt}
\begin{enumerate}[label=\textbf{\alph*)}, leftmargin=1.2cm, itemsep=2pt, topsep=2pt]
  \item \textnormal{[Thông hiểu]} Tổng số tập hợp con của tập hợp$S$là$32$.
  \item \textnormal{[Thông hiểu]} Số tập hợp con của$S$chứa đúng$3$phần tử là$10$.
  \item \textnormal{[Vận dụng thấp]} Số tập hợp con của$S$có tổng các phần tử là một số chẵn bằng$16$.
  \item \textnormal{[Vận dụng thấp]} Có đúng$6$tập hợp con của$S$chứa phần tử$1$nhưng không chứa phần tử$5$.
\end{enumerate}
\vspace{6pt}


% ==============================================================================
% PHẦN III: TRẮC NGHIỆM TRẢ LỜI NGẮN (2,0 ĐIỂM)
% ==============================================================================
\vspace{6pt}
\noindent
\textbf{\large PHẦN III. Câu trắc nghiệm trả lời ngắn (2,0 điểm)}\\
\textit{\small Thí sinh trả lời từ câu 1 đến câu 4. Mỗi câu trả lời đúng thí sinh được 0,5 điểm.}
\vspace{6pt}
\noindent\textbf{Câu 1.} [Thông hiểu] Cho tập hợp$A = \{x \in \mathbb{N}^* \mid 100 \le x \le 999 \text{ và tổng các chữ số của } x \text{ bằng } 3\}$. Hỏi tập hợp$A$có bao nhiêu phần tử?\par\vspace{2pt}
\noindent\textit{Đáp số:} \dotfill\par\vspace{6pt}

\noindent\textbf{Câu 2.} [Thông hiểu] Cho tập hợp$M = \{x \in \mathbb{N} \mid \frac{2x + 15}{x + 1} \in \mathbb{N}\}$. Tập hợp$M$có bao nhiêu phần tử?\par\vspace{2pt}
\noindent\textit{Đáp số:} \dotfill\par\vspace{6pt}

\noindent\textbf{Câu 3.} [Vận dụng cao] Một lớp học có 42 học sinh. Trong đó có 25 bạn tham gia câu lạc bộ Bóng đá, 20 bạn tham gia câu lạc bộ Cờ vua, và 6 bạn không tham gia câu lạc bộ nào trong hai câu lạc bộ trên. Hỏi có bao nhiêu bạn tham gia cả hai câu lạc bộ?\par\vspace{2pt}
\noindent\textit{Đáp số:} \dotfill\par\vspace{6pt}

\noindent\textbf{Câu 4.} [Vận dụng cao] Tìm số tự nhiên$n$nhỏ nhất sao cho tập hợp$A = \{n + 1; n + 2; n + 3; n + 4; n + 5\}$không chứa bất kỳ số nguyên tố nào.\par\vspace{2pt}
\noindent\textit{Đáp số:} \dotfill\par\vspace{6pt}


% ==============================================================================
% PHẦN IV: TỰ LUẬN (3,0 ĐIỂM)
% ==============================================================================
\vspace{6pt}
\noindent
\textbf{\large PHẦN IV. Tự luận (3,0 điểm)}\\
\textit{\small Thí sinh trình bày chi tiết các bước giải từ câu 1 đến câu 3 vào giấy làm bài.}
\vspace{6pt}
\noindent\textbf{Bài 1 (1,0 điểm).} [Thông hiểu] Cho hai tập hợp$A = \{2x + 1 \mid x \in \mathbb{N},\, x \le 4\}$và$B = \{y \in \mathbb{N} \mid (y - 3)(y - 5)(y^2 - 49) = 0\}$.\par\noindent a) Viết tập hợp$A$và tập hợp$B$bằng cách liệt kê các phần tử.\par\noindent b) Xác định tập hợp$C$gồm các phần tử vừa thuộc$A$vừa thuộc$B$, và tìm tất cả các tập con của$C$.\par\vspace{8pt}

\noindent\textbf{Bài 2 (1,0 điểm).} [Vận dụng cao] Trong đợt khảo sát sở thích đọc sách của 50 học sinh khối 6, người ta thu được dữ liệu như sau:\par\noindent - Có 30 bạn thích đọc sách Khoa học;\par\noindent - Có 25 bạn thích đọc sách Văn học;\par\noindent - Có 10 bạn không thích cả hai loại sách trên.\par\noindent Gọi$x$là số bạn chỉ thích đúng một trong hai loại sách trên. Hãy sử dụng kiến thức về tập hợp để tìm giá trị của$x$.\par\vspace{8pt}

\noindent\textbf{Bài 3 (1,0 điểm).} [Vận dụng cao] Cho tập hợp$S = \{1; 2; 3; \dots; 20\}$. Tìm số phần tử lớn nhất của một tập con$A \subset S$sao cho trong tập hợp$A$không có hai phần tử nào có tổng bằng$21$.\par\vspace{8pt}


\begin{center}
\textbf{---------- HẾT ----------}\\
\textit{\small Cán bộ coi thi không giải thích gì thêm.}
\end{center}

\newpage
% ==============================================================================
% HƯỚNG DẪN CHẤM VÀ ĐÁP ÁN CHI TIẾT
% ==============================================================================
\begin{center}
{\Large\bfseries HƯỚNG DẪN CHẤM VÀ ĐÁP ÁN CHI TIẾT}\\[4pt]
{\large\bfseries Toán 6 (KNTT) - Bài 1: Tập hợp - Bài kiểm tra số 3}\\[2pt]
{\bfseries Mã đề: \texttt{c1\_6\_bai1\_test3}}
\end{center}

\vspace{10pt}
\noindent\textbf{\large 1. ĐÁP ÁN PHẦN I: TRẮC NGHIỆM NHIỀU PHƯƠNG ÁN (3,0 điểm)}
\vspace{4pt}
\begin{center}
\begin{tabular}{|c|c|c|c|c|c|c|c|c|c|c|c|}
\hline
\textbf{1} & \textbf{2} & \textbf{3} & \textbf{4} & \textbf{5} & \textbf{6} & \textbf{7} & \textbf{8} & \textbf{9} & \textbf{10} & \textbf{11} & \textbf{12} \\
\hline
\textbf{A} & \textbf{A} & \textbf{C} & \textbf{B} & \textbf{A} & \textbf{A} & \textbf{B} & \textbf{B} & \textbf{A} & \textbf{A} & \textbf{A} & \textbf{A} \\
\hline
\end{tabular}
\end{center}
\vspace{6pt}


\noindent\textbf{\large 2. ĐÁP ÁN PHẦN II: TRẮC NGHIỆM ĐÚNG / SAI (2,0 điểm)}
\vspace{4pt}
\begin{center}
\begin{tabular}{|c|c|c|c|c|}
\hline
\textbf{Câu hỏi} & \textbf{Ý a)} & \textbf{Ý b)} & \textbf{Ý c)} & \textbf{Ý d)} \\
\hline
\textbf{Câu 1} & \textbf{Đúng} & \textbf{Sai} & \textbf{Sai} & \textbf{Đúng} \\
\hline
\textbf{Câu 2} & \textbf{Đúng} & \textbf{Đúng} & \textbf{Đúng} & \textbf{Sai} \\
\hline
\end{tabular}
\end{center}
\vspace{6pt}


\noindent\textbf{\large 3. ĐÁP ÁN PHẦN III: TRẢ LỜI NGẮN (2,0 điểm)}
\vspace{4pt}
\begin{center}
\begin{tabular}{|c|m{11cm}|c|}
\hline
\textbf{Câu} & \textbf{Tóm tắt yêu cầu} & \textbf{Đáp số chính xác} \\
\hline
\textbf{Câu 1} & Cho tập hợp $A = \{x \in \mathbb{N}^* \mid 100 \le x \le 999 \text{ và tổng các... & \textbf{6} \\
\hline
\textbf{Câu 2} & Cho tập hợp$M = \{x \in \mathbb{N} \mid \frac{2x + 15}{x + 1} \in \mathbb{N}\}$... & \textbf{2} \\
\hline
\textbf{Câu 3} & Một lớp học có 42 học sinh. Trong đó có 25 bạn tham gia câu lạc bộ Bóng đá, 20 b... & \textbf{9} \\
\hline
\textbf{Câu 4} & Tìm số tự nhiên$n$nhỏ nhất sao cho tập hợp $A = \{n + 1; n + 2; n + 3; n + 4;... & \textbf{23} \\
\hline
\end{tabular}
\end{center}
\vspace{6pt}


\noindent\textbf{\large 4. LỜI GIẢI CHI TIẾT VÀ BIỂU ĐIỂM TỰ LUẬN (3,0 điểm)}
\vspace{4pt}
\noindent\textbf{\underline{Bài 1 (1,0 điểm):}}\par\vspace{3pt}
Bước 1 (0.5đ): Liệt kê đúng$A = \{1; 3; 5; 7; 9\}$. Với$B$, giải phương trình$(y-3)(y-5)(y^2-49)=0$trong$\mathbb{N}$được$y \in \{3; 5; 7\}$(loại$y = -7$vì$y \in \mathbb{N}$), do đó$B = \{3; 5; 7\}$.\par\noindent Bước 2 (0.5đ): Xác định tập hợp$C = \{3; 5; 7\}$. Liệt kê đúng đủ 8 tập hợp con của$C$:$\emptyset$,$\{3\}$,$\{5\}$,$\{7\}$,$\{3; 5\}$,$\{3; 7\}$,$\{5; 7\}$,$\{3; 5; 7\}$.\par\vspace{8pt}

\noindent\textbf{\underline{Bài 2 (1,0 điểm):}}\par\vspace{3pt}
Bước 1 (0.5đ): Gọi$K$là tập các bạn thích Khoa học,$V$là tập các bạn thích Văn học. Số học sinh thích ít nhất một trong hai loại sách là:$50 - 10 = 40$(bạn).\par\noindent Bước 2 (0.5đ): Số học sinh thích cả hai loại sách là:$|K \cap V| = |K| + |V| - |K \cup V| = 30 + 25 - 40 = 15$(bạn). Số học sinh chỉ thích đúng một trong hai loại sách là:$x = (30 - 15) + (25 - 15) = 15 + 10 = 25$(bạn).\par\vspace{8pt}

\noindent\textbf{\underline{Bài 3 (1,0 điểm):}}\par\vspace{3pt}
Bước 1 (0.5đ): Chia tập hợp$S$thành 10 cặp số có tổng bằng 21:$\{1; 20\}$,$\{2; 19\}$,$\{3; 18\}$,$\dots$,$\{10; 11\}$. Để không có hai phần tử nào có tổng bằng 21, từ mỗi cặp số trên ta chỉ được chọn nhiều nhất một số vào tập$A$.\par\noindent Bước 2 (0.5đ): Vì có 10 cặp rời nhau nên số phần tử của$A$không vượt quá 10, tức là$|A| \le 10$. Chọn ví dụ tập$A = \{1; 2; 3; 4; 5; 6; 7; 8; 9; 10\}$thỏa mãn điều kiện vì tổng của hai số bất kỳ trong$A$luôn nhỏ hơn hoặc bằng$9 + 10 = 19 < 21$. Vậy số phần tử lớn nhất cần tìm là 10.\par\vspace{8pt}

\end{document}
